\documentclass[12pt,a4paper]{memoir}

% Basic setup for a handout
\usepackage[T1]{fontenc}
\usepackage{newpxtext,newpxmath}
% To enable compactitem
\usepackage{paralist}
\usepackage{tabularx}
\usepackage{enumitem}
\usepackage{nomencl} % For \nomenclature
\usepackage{makeidx} % For \index
\usepackage{ragged2e} % For \RaggedRight
\makeindex
\makenomenclature
\setsecnumdepth{subsection} % Number sections and subsections

% Allow cross referencing to another document
\usepackage{xr}
\externaldocument{ada_volume1}


% Allows background color to paragraphs
\usepackage[framemethod=tikz]{mdframed} 

% Automatically numbered Exercises
\newcounter{exercise}[chapter]
\renewcommand{\theexercise}{\thechapter.\arabic{exercise}}
\newenvironment{exercise}
  {\refstepcounter{exercise}\par\medskip\noindent\textbf{\theexercise}\quad}
  {\par\medskip}

% Page layout
\setlength{\textwidth}{6.5in}
\setlength{\textheight}{9in}
\setlength{\oddsidemargin}{0in}
\setlength{\evensidemargin}{0in}
\setlength{\topmargin}{0in}
\setlength{\headheight}{0.5in}
\setlength{\headsep}{0.25in}

% To embed flowcharts
\usepackage{tikz} 
  \usetikzlibrary{shapes, arrows, positioning, calc, math, arrows.meta}
% Define block styles
\tikzstyle{decision} = [diamond, draw, fill=gray, 
    text width=4.5 em, text badly centered, node distance=3 cm, inner sep = 0 pt]
\tikzstyle{whiteblock} = [rectangle, draw, node distance=2.5 cm, text width=7.2 em, text centered, rounded corners, minimum height=4 em]
\tikzstyle{widewhiteblock} = [rectangle, draw, node distance=2.5 cm, text width=4cm, text centered, rounded corners, minimum height=4 em]
\tikzstyle{blueblock} = [rectangle, draw, fill=gray!20, text width=7.2 em, text centered, rounded corners, minimum height=4 em, node distance=3 cm]  
\tikzstyle{wideblueblock} = [rectangle, draw, fill=gray!20, 
    text width=4cm, text centered, rounded corners, minimum height=4 em,
    node distance=3 cm]  
\tikzstyle{line} = [draw, -latex']
\tikzstyle{cloud} = [draw, circle, fill=red!20, node distance=3 cm,
    minimum height=2 em]


\begin{document}

% Title page
\thispagestyle{empty}
\begin{center}
    \vspace*{2in}
    {\Huge \textbf{Audit Data Analysis -- Volume 1}} \\[0.3in]
    {\LARGE \textit{Review Questions}} \\[0.5in]
    {\Large Lucas Hoogduin} \\[0.5in]
    {\large \today}
\end{center}
\newpage

% Table of contents (optional for handout)
\tableofcontents
\newpage

% Main content

\chapter{Probability distributions}

\chapter{Estimating the population mean and proportion}

\chapter{Estimation with auxiliary variables and stratification}

\chapter{Hypothesis testing}

\chapter{Regression analysis}

\chapter{Goodness of fit}

\begin{exercise}
Explain in your own words what it means for the observed frequencies to be close to the expected frequencies.
\end{exercise}

\begin{exercise}
    Explain why expected frequencies do not have to be whole numbers.
\end{exercise}

\begin{exercise}
    Use the customer satisfaction example from Section \ref{sec:the_chi_square_goodness_of_fit_test}. 
    \begin{enumerate}
        \item Write down the null hypothesis in words. 
        \item Write down the alternative hypothesis in words.
    \end{enumerate}
    
\end{exercise}
    
\begin{exercise}
    Explain why Benford's Law does not imply that all first digits are equally likely.
\end{exercise}

\begin{exercise}
    Refer to the Inventory first-digit data in Section \ref{sec:the_chi_square_distribution}. Explain why rejection of the null hypothesis does not prove fraud.
\end{exercise}

\begin{exercise}
    An advantage of the chi-square statistic is that its individual components can be examined separately. Why is this useful in an audit setting?
\end{exercise}

\begin{exercise}
    Suppose an auditor is considering whether Benford's Law is appropriate for each of the following datasets:

\begin{itemize}
    \item invoice amounts from a large supplier ledger;
    \item employee identification numbers;
    \item retail prices ending in 95 or 99 cents;
    \item total annual sales by customer;
    \item postal codes.
\end{itemize}

For each dataset:

\begin{enumerate}
    \item State whether Benford's Law is likely to be appropriate.
    \item Give one reason for your answer.
    \item Identify whether the data are naturally occurring measurements, assigned numbers, or values constrained by human convention.
\end{enumerate}
\end{exercise}

\begin{exercise}
    You are analyzing second- and third-digit frequencies from a large dataset. Explain why the digit zero can occur in this analysis.
\end{exercise}

\begin{exercise}
    In a chi-square goodness-of-fit test, the expected frequency in each category should generally be sufficiently large.
    \begin{enumerate}
    \item For a discrete uniform distribution with ten categories, calculate the
    minimum sample size needed to obtain \(E_i \geq 5\) in every category.
    \item Repeat the calculation for a first-digit Benford test with nine
    categories. Use the smallest Benford probability in your calculation.
    \item Explain why first-two-digit Benford tests require much larger samples
    than first-digit tests.
\end{enumerate}
    
\end{exercise}

\begin{exercise}
    Explain why the chi-square goodness-of-fit test is not the preferred test for normality in this small-sample example.
\end{exercise}

\begin{exercise}
Assume that you are part of an audit team and have received a dataset containing
transaction amounts from a client.

Design a simple audit analytics procedure.

\begin{enumerate}
    \item State the audit question you want to investigate.
    \item Decide whether Benford's Law, a discrete uniform distribution, or
    another expected distribution is appropriate.
    \item Explain why the chosen distribution is appropriate for the data.
    \item  Write a short paragraph explaining the result to a non-statistical audit manager.
\end{enumerate}

Assume that the number 7 occurs more frequently than expected. Explain how you would investigate this anomaly.

\begin{enumerate}[resume]
    \item Explain why this result is not sufficient evidence of fraud.
    \item Write a short audit conclusion.
    \item Recommend one additional audit procedure based on the results.
\end{enumerate}

At the same time, you observe that the number 2 occurs less frequently than expected. 
\begin{enumerate}[resume]
    \item Explain how you would investigate this anomaly.
\end{enumerate}


\end{exercise}

\end{document}